% 自动生成：python3 scripts/generate-code.py，以 main-landscape.tex 为模板；请修改横版入口或章节源文件后重新生成。
\documentclass{article}

% 在主文件的 \documentclass 之后加载；建议使用 XeLaTeX 编译。
% 此文件集中加载宏包，具体版式与自定义命令分别放在 layout-portrait.tex 和 mystyle.tex 中。

% 中文支持；不采用 ctex 的中文标题名称与版式方案。
% 无布局的独立接口测试使用 Fandol；正式字体由 layout-portrait.tex 覆盖。
\PassOptionsToPackage{fontset=fandol}{ctex}
\usepackage[scheme=plain]{ctex}

% 数学公式、矩阵、对齐环境与数学符号（含 \mathbb）。
\usepackage{amsmath}
\usepackage{amssymb}
\usepackage{unicode-math}

% 图片与颜色。
\usepackage{graphicx}
\usepackage{xcolor}

% 删除线；normalem 保留 \emph 的原有强调语义。
\usepackage[normalem]{ulem}

% 页面尺寸、边距与多栏排版；具体参数在 layout-portrait.tex 中设置。
\usepackage{geometry}
\usepackage{multicol}
\usepackage{titlesec}
\usepackage{needspace}
\usepackage{enumitem}
\usepackage{fancyhdr}
\usepackage{tocloft}
\usepackage{lastpage}
\usepackage{etoolbox}
\usepackage{newunicodechar}

% 代码高亮，以及后续代码标题栏、引用块所需的盒子支持。
\usepackage{minted}
\usepackage{tcolorbox}
\tcbuselibrary{skins,breakable,minted}

% 目录链接、URL 和交叉引用；在上述宏包之后加载。
\usepackage[
    hidelinks
]{hyperref}

% 代码版复用横版布局；专用覆盖写在载入之后。
% 横版布局：只放本版参数；三版共用的实现及双栏机制见 preamble/shared/。
% A4 横向双栏；四边 12 mm，上下包含页眉、页脚及其间距。
\geometry{
  a4paper,landscape,left=12mm,right=12mm,top=12mm,bottom=12mm,
  includeheadfoot,headheight=12pt,headsep=4mm,footskip=6mm
}

% 正文、数学、行内等宽与代码块字体及单色符号见共用文件。
% 三版共用的字体：正文与导航、单色符号、数学、行内等宽与代码块。
% 由各版布局在 \geometry 之后 \input；字号、行距等参数由各版布局设置。
% 本机字体须可供 XeLaTeX 访问，不静默替换。

% 正文与导航字体：使用与已确认样张一致的 Latin Modern Roman 10 pt 静态字形。
\setmainfont{lmroman10}[
  Extension=.otf,UprightFont=*-regular,BoldFont=*-bold,
  ItalicFont=*-italic,BoldItalicFont=*-bolditalic
]
\setsansfont{PingFang SC}[BoldFont={PingFang SC Semibold}]
\setCJKmainfont{Songti SC}[BoldFont={Songti SC Bold},ItalicFont={Songti SC}]
\setCJKsansfont{PingFang SC}[BoldFont={PingFang SC Semibold}]
% 宋体正文的双引号改用当前西文正文字形（用户确认），输入仍为 Unicode “”。
% 跟随 NFSS 字体族选择：粗体/斜体保留规则，苹方与 Maple（含自定义代码字体）恢复原分类。
% 保存 xeCJK 已分配的类号，不硬编码；半角左右标点类取自西文圆括号。
% 不在 selectfont 中反复调用 xeCJKDeclareCharClass：它会全局追加字符范围表。
% 这里只作局部 XeTeXcharclass 赋值，嵌套字体分组结束时自然恢复外层分类。
\makeatletter
\chardef\BookQuoteFullLeft=\XeTeXcharclass"201C
\chardef\BookQuoteFullRight=\XeTeXcharclass"201D
\chardef\BookQuoteHalfLeft=\XeTeXcharclass`\(
\chardef\BookQuoteHalfRight=\XeTeXcharclass`\)
\newcommand{\BookSelectQuoteClass}{%
  \ifx\f@family\rmdefault
    \XeTeXcharclass"201C=\BookQuoteHalfLeft
    \XeTeXcharclass"201D=\BookQuoteHalfRight
  \else
    \XeTeXcharclass"201C=\BookQuoteFullLeft
    \XeTeXcharclass"201D=\BookQuoteFullRight
  \fi}
\AddToHook{selectfont}[book/quotes]{\BookSelectQuoteClass}
\makeatother
% 单色符号使用项目内的静态字体；标题中也保持常规直立字重。
% 只映射原稿实际使用的两个字符，不更改正文或全局字符分类。
\newfontfamily\BookSymbolFont{NotoEmoji-Regular.ttf}[
  Path=preamble/fonts/,
  BoldFont=NotoEmoji-Regular.ttf,
  ItalicFont=NotoEmoji-Regular.ttf,
  BoldItalicFont=NotoEmoji-Regular.ttf
]
% xeCJK 默认借用当前字体空格；Noto Emoji 的空格宽达 1.27em。
% 仅在符号局部分组内固定中西文间距，避免中文后的符号被推远。
\DeclareRobustCommand{\BookStar}{{%
  \xeCJKsetup{CJKecglue={\hskip .25em\relax}}%
  \BookSymbolFont\char"2B50\relax}}
\DeclareRobustCommand{\BookFlag}{{%
  \xeCJKsetup{CJKecglue={\hskip .25em\relax}}%
  \BookSymbolFont\char"1F6A9\relax}}
\newunicodechar{⭐}{\BookStar}
\newunicodechar{🚩}{\BookFlag}
% 书签保留原 Unicode，避免把字体命令写入 PDF 字符串。
\pdfstringdefDisableCommands{%
  \def⭐{\string⭐}%
  \def🚩{\string🚩}%
}
\setmathfont{Latin Modern Math}
% 直立数学字母与算子均使用数学字体；\text 仍使用正文字体。
% 显式外层分组兼容原稿的 ^\mathrm{mid}、_\mathrm{mid} 写法。
% 直接使用 mathrm=sym 时，unicode-math 的内部组命令不能充当上下标组。
% 在 unicode-math 完成文档起始初始化后设置，避免被其默认定义覆盖。
\AddToHook{begindocument/end}{\renewcommand{\mathrm}[1]{{\symup{#1}}}}
\setoperatorfont{\symup}
% 保留 \boldsymbol 对完整表达式（含算子和定界符）的语义。
% 数学字体没有独立粗体文件，bold 版本采用轻度合成加粗。
\setmathfont{Latin Modern Math}[version=bold,FakeBold=1.2]
% 保留文本类数学字母（如 \mathbf、\mathit）的 bold 版本字体映射。
% \mathrm 已由上方定义改用数学字形，不走此文本映射。
\setboldmathrm{lmroman10}[
  Extension=.otf,UprightFont=*-bold,BoldFont=*-bold,
  ItalicFont=*-bolditalic,BoldItalicFont=*-bolditalic
]

% 行内等宽文字使用 Maple，字号随所在正文；与代码块采用相同字形特性。
\setmonofont{Maple Mono NF CN}[
  Ligatures={NoCommon,NoContextual,NoDiscretionary,NoHistoric},
  RawFeature={-calt;-tlig;+cv01;+cv32;+cv33;+cv34;+cv35;+cv36;+cv37}
]
\setCJKmonofont{Maple Mono NF CN}[
  Ligatures={NoCommon,NoContextual,NoDiscretionary,NoHistoric},
  RawFeature={-calt;-tlig;+cv01;+cv32;+cv33;+cv34;+cv35;+cv36;+cv37}
]

% 只供 minted 使用（fontfamily=bookcode）；不改变 \texttt 或代码标题栏。
% 禁用标准、上下文、可选、历史连字，以及 calt/TeX 字符替换。
% cv01 去除符号间隙；cv32--cv37 将手写斜体改为常规斜体（仍保留倾斜）。
\newfontfamily\BookCodeLatin{Maple Mono NF CN}[
  Ligatures={NoCommon,NoContextual,NoDiscretionary,NoHistoric},
  RawFeature={-calt;-tlig;+cv01;+cv32;+cv33;+cv34;+cv35;+cv36;+cv37}
]
\newCJKfontfamily\BookCodeCJK{Maple Mono NF CN}[
  Ligatures={NoCommon,NoContextual,NoDiscretionary,NoHistoric},
  RawFeature={-calt;-tlig;+cv01;+cv32;+cv33;+cv34;+cv35;+cv36;+cv37}
]
\makeatletter
\@namedef{FV@fontfamily@bookcode}{%
  \def\FV@FontScanPrep{}%
  \def\FV@FontFamily{\BookCodeLatin\BookCodeCJK}%
}
\makeatother

% 页眉页脚、目录点线与代码配色见共用文件。
% 三版共用的配色，由各版布局在 shared/fonts.tex 之后 \input。
% 页眉页脚、目录、代码样式与扉页引用这里的颜色；mystyle.tex 中各接口自带的颜色不在此处。

% 页眉、页脚与次要文字。
\definecolor{BookMuted}{HTML}{666666}
% 目录二、三级条目的点线。
\definecolor{BookTOCDots}{HTML}{888888}
% 代码阅读 / 打印配色：仅两种彩色；名称、调用与标点保持深灰。
\definecolor{BookCodeText}{HTML}{30343B}
\definecolor{BookCodeBlue}{HTML}{506B82}
\definecolor{BookCodeGreen}{HTML}{58715D}
\definecolor{BookCodeComment}{HTML}{70757D}
% 扉页：主要文字、次要文字与细框；唯一强调色与代码高亮的蓝灰一致。
\definecolor{BookTitlePageInk}{HTML}{1D1D1F}
\definecolor{BookTitlePageSecondary}{HTML}{6E6E73}
\definecolor{BookTitlePageHairline}{HTML}{D1D1D6}
\colorlet{BookTitlePageAccent}{BookCodeBlue}

\makeatletter
% 固定正文及常用相对字号，数学上下间距随正常字号一起恢复。
\renewcommand\normalsize{%
  \@setfontsize\normalsize{10pt}{14pt}%
  \abovedisplayskip=7pt plus 2pt minus 2pt
  \belowdisplayskip=7pt plus 2pt minus 2pt
  \abovedisplayshortskip=3pt plus 1pt
  \belowdisplayshortskip=5pt plus 1pt minus 1pt
}
\renewcommand\small{\@setfontsize\small{9pt}{12pt}}
\renewcommand\footnotesize{\@setfontsize\footnotesize{8pt}{10.5pt}}
\makeatother
\normalsize
\setlength{\parindent}{0pt}
\setlength{\parskip}{2pt plus 1pt minus 1pt}
% 窄栏含较长行内公式时，增加断行算法的应急伸展量；不改变正常行距。
% 用户确认 6em（曾试改 3em 后撤回）：3em 时 Overfull 较多，6em 只留少量过松行。
\setlength{\emergencystretch}{6em}
\raggedbottom
\clubpenalty=10000
\widowpenalty=10000
\displaywidowpenalty=10000
\setlist{topsep=3pt,itemsep=1pt,parsep=1pt,partopsep=0pt,leftmargin=*}
% 二级无序列表用空心圆点（用户确认）；Latin Modern 的 U+25E6 与一级圆点外框相同，只差填充。
\renewcommand{\labelitemii}{\labelitemfont\textopenbullet}

\titleformat{\section}{\sffamily\bfseries\fontsize{16}{21}\selectfont}{}{0pt}{\BookChapterAnchor}
\titleformat{\subsection}{\sffamily\bfseries\fontsize{12}{17}\selectfont}{}{0pt}{}
\titleformat{\subsubsection}{\sffamily\bfseries\fontsize{10.5}{15}\selectfont}{}{0pt}{}
\titlespacing*{\section}{0pt}{10pt}{5pt}
\titlespacing*{\subsection}{0pt}{8pt plus 2pt minus 2pt}{4pt}
\titlespacing*{\subsubsection}{0pt}{6pt plus 2pt minus 1pt}{3pt}
\titleformat{\paragraph}{\sffamily\bfseries\fontsize{10}{14}\selectfont}{}{0pt}{}
\titleformat{\subparagraph}{\sffamily\bfseries\fontsize{10}{14}\selectfont}{}{0pt}{}
\titlespacing*{\paragraph}{0pt}{5pt}{2pt}
\titlespacing*{\subparagraph}{0pt}{4pt}{2pt}
% 栏距；启动双栏、通栏目录标题、页顶所在章的页眉与图片按栏宽见 shared/twocolumn.tex。
\setlength{\columnsep}{8mm}

% 标题编号、目录条目与书签设置，以及双栏机制见共用文件。
% 三版共用的标题编号、目录条目与书签设置，由各版布局在设置字号之后 \input。
% 目录标题及目录前后的分页由各版负责：竖版用 tocloft 的标题设置，双栏版本见 shared/twocolumn.tex。
% 每次排目录前须把 \BookTOCPreviousLevel 清零：竖版在布局中挂钩，双栏版本在重定义的 \tableofcontents 中清零。

% 保留原文的 0xPQ 编号：只关闭章节标题的自动编号。
\setcounter{secnumdepth}{-1}
% 目录展开到三级；PDF 书签包含全部标题层级并默认全部展开。
\setcounter{tocdepth}{3}
\hypersetup{bookmarksdepth=5,bookmarksopen=true}
\renewcommand{\contentsname}{目录}
% 方案 A：一级沿用苹方 Semibold；二三级沿用正文宋体与 Latin Modern Roman。
\newlength{\BookTOCUnit}
\settowidth{\BookTOCUnit}{\normalfont\fontsize{10}{15}\selectfont 中}
% 每条目录装入完整段落盒：续行保持本级缩进，页码固定在末行右端。
% 条目间使用可丢弃的竖直间距，页首不残留标题前留白。
% 根据前后层级禁止父标题与首个子条目之间分页，不捆绑整个章节。
\makeatletter
\newcount\BookTOCPreviousLevel
\BookTOCPreviousLevel=0
\newcommand{\BookTOCEntry}[4]{%
  \par
  \ifnum\BookTOCPreviousLevel=1 \penalty10000
  \else\ifnum\BookTOCPreviousLevel=2
    \ifnum#1=3 \penalty10000\else\penalty0\fi
  \else\penalty0\fi\fi
  \begingroup
  \@tempdima=0pt
  \ifnum\BookTOCPreviousLevel=1 \advance\@tempdima by5pt\fi
  \ifnum#1=1 \advance\@tempdima by12pt\fi
  \vskip\@tempdima
  \nointerlineskip
  \vbox{%
    \hsize=\linewidth
    \normalfont\fontsize{10}{15}\selectfont
    \ifnum#1=1 \sffamily\bfseries\fontsize{12}{15}\selectfont
    \else\ifnum#1=2 \fontsize{10.5}{15}\selectfont\fi\fi
    \leftskip=#2\relax \rightskip=30pt
    \parfillskip=-\rightskip \parindent=0pt \parskip=0pt
    \interlinepenalty=10000
    \noindent\strut#3\nobreak
    \ifnum#1=1 \hfill
    % 9pt 是实际点中心间距；两端各留 3pt，空间不足一格时不绘点。
    \else{\normalfont\fontsize{10}{15}\selectfont\color{BookTOCDots}%
      \kern3pt\leaders\hbox to9pt{\hss.\hss}\hfill\kern3pt}\fi
    \nobreak\hb@xt@20pt{\hfil\normalfont
      \ifnum#1=1 \fontsize{11}{15}\else\fontsize{10}{15}\fi\selectfont#4}\strut\par
  }%
  \endgroup
  \global\BookTOCPreviousLevel=#1\relax}
\renewcommand*{\l@section}[2]{%
  \ifnum\c@tocdepth>0 \BookTOCEntry{1}{0pt}{#1}{#2}\fi}
\renewcommand*{\l@subsection}[2]{%
  \ifnum\c@tocdepth>1 \BookTOCEntry{2}{2\BookTOCUnit}{#1}{#2}\fi}
\renewcommand*{\l@subsubsection}[2]{%
  \ifnum\c@tocdepth>2 \BookTOCEntry{3}{4\BookTOCUnit}{#1}{#2}\fi}
\makeatother

% 横版与代码版共用的双栏机制：启动双栏、通栏目录标题、页顶所在章的页眉、图片按栏宽解释 \textwidth。
% 由两版布局在 shared/toc.tex 之后 \input；栏距等版面参数由各版布局设置，
% 一级标题格式须在标题文字处调用 \BookChapterAnchor。

% 入口无排版选项；启动标准双栏输出，章节自然连续，不平衡章末。
\setlength{\columnseprule}{0pt}
\AddToHook{begindocument/end}{\twocolumn\raggedbottom}

% 目录标题由下方重写的 \tableofcontents 以通栏居中排出（22/28 pt，后接 16pt），
% 不使用 tocloft 的标题设置。
\makeatletter
% tocloft 在 begindocument 中重定义入口，此处在其初始化结束后安装。
% 标准双栏的可选通栏头只用于目录首张；条目仍按栏自然流动。
\AddToHook{begindocument/end}{%
  \renewcommand{\tableofcontents}{%
    \clearpage
    \markboth{}{}%
    \global\BookTOCPreviousLevel=0
    \twocolumn[{%
      \centering\sffamily\bfseries\fontsize{22}{28}\selectfont
      \contentsname\par\vspace{16pt}%
    }]%
    \@starttoc{toc}%
    \clearpage
  }%
}
\makeatother

% 页眉显示页顶所在的一级标题（用户确认，2026-10-14）：
% 本页没有新章，或新章不在左栏顶端时，取页首延续的章（TopMark）；
% 本页第一个新章恰好位于左栏顶端时，取该章。章节连续排版，
% 一页可含多章，不能沿用竖版取本页最后标记的 \leftmark。
% 位置由标题文字内的 \pdfsavepos 记录，需两轮编译；run.sh 至少两轮。
\makeatletter
\NewMarkClass{book-chapter}
\newcounter{BookChapter}
\AddToHook{cmd/section/before}{\stepcounter{BookChapter}}
% 置于标题文字内部，保证位置与标题同在一栏，不会因断页与标题分离。
\newcommand{\BookChapterAnchor}[1]{%
  \leavevmode\pdfsavepos
  \RecordProperties{BookChapter@\the\value{BookChapter}}{xpos,ypos}#1}
\newdimen\BookChapterTopLimit
\newcommand{\BookHeaderChapter}{%
  \begingroup
  \edef\BookTop{\TopMark{book-chapter}}%
  \edef\BookFirst{\FirstMark{book-chapter}}%
  \let\BookShown\BookTop
  \ifx\BookTop\BookFirst\else
    \ifx\BookTop\@empty\let\BookShown\BookFirst\fi
    % 左栏顶端标题的基线距版心顶约一个 topskip；留 20pt 容差，
    % 上方哪怕只有一行正文，标题基线也会再低约 30pt。
    \BookChapterTopLimit=\dimexpr\paperheight-1in-\voffset-\topmargin
      -\headheight-\headsep-\topskip-20pt\relax
    \ifdim\RefProperty{BookChapter@\BookFirst}{xpos}sp<.5\paperwidth
      \ifdim\RefProperty{BookChapter@\BookFirst}{ypos}sp>\BookChapterTopLimit
        \let\BookShown\BookFirst
      \fi
    \fi
  \fi
  \ifx\BookShown\@empty\else
    \csname BookChapterName@\BookShown\endcsname
  \fi
  \endgroup}
\makeatother
\pagestyle{fancy}
\fancyhf{}
\fancyhead[L]{\sffamily\fontsize{8}{10}\selectfont\color{BookMuted}%
  \BookHeaderChapter}
\fancyfoot[C]{\sffamily\fontsize{8}{10}\selectfont\color{BookMuted}\thepage}
\renewcommand{\headrulewidth}{0pt}
\renewcommand{\footrulewidth}{0pt}
% 保留标准标记，另记章号及章名供页眉判断页顶所在的章。
\renewcommand{\sectionmark}[1]{%
  \markboth{#1}{}%
  \expandafter\gdef\csname BookChapterName@\the\value{BookChapter}\endcsname{#1}%
  \InsertMark{book-chapter}{\the\value{BookChapter}}}
\renewcommand{\subsectionmark}[1]{\markright{#1}}
\fancypagestyle{plain}{\fancyhf{}\fancyfoot[C]{\sffamily\fontsize{8}{10}\selectfont\color{BookMuted}\thepage}}

% 双栏图片按当前容器计算正文宽度；保持原路径、可选参数和长宽比。
% 用户已批准将现有 0.8\textwidth / \textwidth 分别解释为 0.8 / 1 栏宽。
% 局部分组避免改变双栏输出的整页 textwidth。
\NewCommandCopy{\BookIncludeGraphics}{\includegraphics}
\RenewDocumentCommand{\includegraphics}{s O{} m}{%
  \begingroup
  \setlength{\textwidth}{\linewidth}%
  \IfBooleanTF{#1}
    {\BookIncludeGraphics*[#2]{#3}}
    {\BookIncludeGraphics[#2]{#3}}%
  \endgroup
}


% 代码块的词法配色、minted 默认设置与代码卡片标题见共用文件。
% 三版共用的代码块样式：Pygments 词法配色、minted 默认设置与代码卡片标题。
% 由各版布局在 shared/fonts.tex、shared/colors.tex 之后 \input。
% 某版需要不同的代码字号时，在 \input 之后用 \setminted{fontsize=…} 覆盖。

% 使用 bw 提供 Pygments 转义与分词接口，在每次加载样式后替换 token 外观。
% 兼容目标：minted 3.x / Pygments 2.x；不安装系统 Python 插件。
% 未列出的 token 不增加颜色、粗体、下划线或错误边框。
\makeatletter
\newcommand{\BookCodeBlueText}[1]{\textcolor{BookCodeBlue}{#1}}
\newcommand{\BookCodeGreenText}[1]{\textcolor{BookCodeGreen}{#1}}
\newcommand{\BookCodeCommentText}[1]{\textcolor{BookCodeComment}{#1}}
\def\BookToken@k{\let\PYG@tc\BookCodeBlueText}
\let\BookToken@m\BookToken@k
\let\BookToken@no\BookToken@k
\let\BookToken@ow\BookToken@k
% 内建名称：C++ 的 true、false、NULL，Bash、Python 的内建命令与函数。
\let\BookToken@nb\BookToken@k
\def\BookToken@s{\let\PYG@tc\BookCodeGreenText}
\def\BookToken@c{\let\PYG@tc\BookCodeCommentText}
\def\BookToken@cp{\let\PYG@tc\BookCodeBlueText\let\PYG@it\relax}
\def\BookToken@cpf{\let\PYG@tc\BookCodeGreenText\let\PYG@it\relax}
% Pygments 2.19、2.20 把 1ll、1llu、1lu 分成浮点数 1l 与名称 l、lu、u（pygments/pygments#1847），
% 把 C++23 的 1z、1uz、1zu 分成整数与名称 z、zu。数字后紧跟只由 l、u、z（含大写）组成的名称时，
% 名称只会是拆开的后缀，改按数字着色；10_km 等自定义后缀以下划线开头，不受影响。上游修复进入 TeX Live 后可删除。
% fvextra 给每个 token 加分组，折行时又以 \FancyVerbBreakStart/Stop 包住内容，故用全局状态，匹配时容许这层包装。
% 每行先测宽再排版，每遍都从 \FancyVerbFormatText 开始，在此清除上一遍留下的状态。
\ExplSyntaxOn
\bool_new:N \g__book_code_after_number_bool
\regex_const:Nn \c__book_code_suffix_regex
  { \A (?: \c{FancyVerbBreakStart} )? [lLuUzZ]{1,3} (?: \c{FancyVerbBreakStop} )? \Z }
\cs_new_protected:Npn \BookCode@Suffix #1#2
  {
    \bool_if:NT \g__book_code_after_number_bool
      {
        \str_if_eq:nnT {#1} {n}
          { \regex_match:NnT \c__book_code_suffix_regex {#2} { \BookToken@m } }
      }
    \str_if_eq:eeTF { \str_range:nnn {#1} { 1 } { 3 } } { l+m }
      { \bool_gset_true:N \g__book_code_after_number_bool }
      { \bool_gset_false:N \g__book_code_after_number_bool }
  }
\cs_new_eq:NN \__book_code_format_text:n \FancyVerbFormatText
\cs_set_protected:Npn \FancyVerbFormatText #1
  {
    \bool_gset_false:N \g__book_code_after_number_bool
    \__book_code_format_text:n {#1}
  }
\ExplSyntaxOff
% Pygments 的 Bash 词法把 <0.in、diff 1.out 中文件名开头的数字单独标为数字（Literal.Number），
% 其后的 .in、.out 是不经 \PYG 的普通文字；数字后紧跟“.”时改按普通文字排。
% fvextra 在每个 token 之后结束分组，故越过这个 \endgroup 再看下一个字符；不是 \endgroup 时照常着色。
\def\BookCode@Number#1{\@ifnextchar\endgroup{\BookCode@Number@i{#1}}{\PYG@do{#1}}}
\def\BookCode@Number@i#1\endgroup{%
  \@ifnextchar.{\let\PYG@tc\relax\PYG@do{#1}\endgroup}{\PYG@do{#1}\endgroup}}
\AddToHook{cmd/minted@defstyle/after}{%
  \edef\BookCodeStyle{\mintedtexoptvalueof{style}}%
  \ifdefstring{\BookCodeStyle}{bw}{%
    \def\PYG@tok#1{\ifcsname BookToken@#1\endcsname
      \csname BookToken@#1\endcsname\fi}%
    % 同 Pygments 的 \PYG，只加入上面两项处理。
    \def\PYG#1#2{\PYG@reset\PYG@toks#1+\relax+\BookCode@Suffix{#1}{#2}%
      \ifstrequal{#1}{l+m}{\BookCode@Number{#2}}{\PYG@do{#2}}}%
  }{}%
}
\makeatother
\setminted{
  fontfamily=bookcode,fontsize=\fontsize{8}{10.5}\selectfont,
  baselinestretch=1,style=bw,linenos=false,
  formatcom=\color{BookCodeText}\mdseries\upshape,
  breaklines=true,breakanywhere=false,tabsize=4
}
% 词法高亮并非 C++ 语义分析；仅补充明确要求的三个标准库类型名称、本书的整数类型别名 i64、i128、u64、u32，
% 以及 GCC 的 __int128_t、__uint128_t（__int128 已是关键字）。这些别名与 GCC 类型名不会用作变量名。
\setminted[cpp]{extrakeywordstype={vector,string,pair,i64,i128,u64,u32,__int128_t,__uint128_t}}
\setminted[c++]{extrakeywordstype={vector,string,pair,i64,i128,u64,u32,__int128_t,__uint128_t}}
\setminted[cxx]{extrakeywordstype={vector,string,pair,i64,i128,u64,u32,__int128_t,__uint128_t}}
\newfontfamily\BookCodeTitleLatin{PingFang SC Medium}
\newCJKfontfamily\BookCodeTitleCJK{PingFang SC Medium}
% coltext 同时保证 tcolorbox 续页恢复深灰，而非默认纯黑。
\tcbset{input-code-title/.style={coltext=BookCodeText,fonttitle=\mdseries\upshape\BookCodeTitleLatin\BookCodeTitleCJK\fontsize{9}{12}\selectfont}}
% 代码卡片与引用块共用的前后距，取代 tcolorbox 默认的 beforeafter skip balanced；
% 长度 \BookBoxBeforeSkip、\BookBoxAfterSkip 由各版布局设定。
% 盒前用 tcolorbox 的 before skip：盒顶距上一行的底为前距，盒与盒之间也是这个距离。
% 紧跟标题时 tcolorbox 不执行盒前设置（\iftcb@nobreak），标题与盒之间不受影响。
% 盒后把 \prevdepth 归零再加后距：后接正文时按段距与行距接排，后接标题时盒底相当于上一行的基线。
\makeatletter
\tcbset{book-box-spacing/.style={
  before skip=\BookBoxBeforeSkip,
  after={\tcb@parfillskip@check\par\ifvmode\prevdepth\z@\addvspace{\BookBoxAfterSkip}\fi}}}
\makeatother


% 扉页：外观与提交信息读取见共用文件；入口提供书名、作者、英文书名与配图路径。
% 版面沿用竖版的元素与间距，按 A4 横向重排：书名基线距页顶 48 mm，提交信息底边距纸边 17 mm；
% 页高只有 210 mm，配图区由竖版的 130 × 110 mm 收紧为 110 × 70 mm。
% 三版共用的扉页：书名、英文书名、作者、配图，以及底部的仓库地址与 Last commit；设了版本标记时悬挂在书名右侧。
% 由各版布局在 shared/fonts.tex、shared/colors.tex 之后 \input，再用 \TitlePageSetup 给出本版的版面参数。
% 书名、作者、英文书名、配图路径、版本标记与仓库地址由入口提供；路径均相对于项目根目录。
% 绝对定位依赖 TikZ 的 remember picture，至少编译两轮（根目录 run.sh 已保证）。

\newfontfamily\BookTitlePageMediumLatin{PingFang SC Medium}
\newCJKfontfamily\BookTitlePageMediumCJK{PingFang SC Medium}
% 提交信息的眉题：小号大写，加宽字距。
\newfontfamily\BookTitlePageEyebrow{PingFang SC Medium}[LetterSpace=14]
% 提交说明原样输出：此字体族不启用 TeX 连字，-- 等不会被替换。
\newfontfamily\BookTitlePagePlainLatin{PingFang SC}

% 版面参数：以纸张中心为原点，书名、英文书名、作者的基线，配图区中心，底部信息区底边的纵向位置，
% 以及配图区的宽和高。配图在配图区内等比缩放后居中；标题区与底部信息区的位置不随图片比例变化。
% 各版布局须给出全部 7 项，缺少时编译报错；数值按原文交给 TikZ 与 \setlength，不预先换算。
\newlength{\BookTitlePageImageWidth}
\newlength{\BookTitlePageImageHeight}
\pgfkeys{
  /book/titlepage/.is family,
  /book/titlepage/title/.store in=\BookTitlePageTitleY,
  /book/titlepage/english/.store in=\BookTitlePageEnglishY,
  /book/titlepage/author/.store in=\BookTitlePageAuthorY,
  /book/titlepage/image/.store in=\BookTitlePageImageY,
  /book/titlepage/stamp/.store in=\BookTitlePageStampY,
  /book/titlepage/image-width/.code={\setlength{\BookTitlePageImageWidth}{#1}},
  /book/titlepage/image-height/.code={\setlength{\BookTitlePageImageHeight}{#1}},
}
\newcommand{\TitlePageSetup}[1]{\pgfkeys{/book/titlepage/.cd,#1}}

\makeatletter
% 元信息：英文书名、配图路径与版本标记；不写或参数为空时省略对应内容。
% 版本标记只写文字（如“代码版”），全角括号与样式由扉页添加。
\newcommand*{\EnglishTitle}[1]{\gdef\BookTitlePage@EnglishTitle{#1}}
\newcommand*{\TitlePageImage}[1]{\gdef\BookTitlePage@Image{#1}}
\newcommand*{\EditionLabel}[1]{\gdef\BookTitlePage@Edition{#1}}
\let\BookTitlePage@EnglishTitle\@empty
\let\BookTitlePage@Image\@empty
\let\BookTitlePage@Edition\@empty
% PDF 文档属性同步入口的书名与作者；有版本标记时标题附上，如“算法竞赛小记（代码版）”。
% 入口未设置书名或作者时不写入。
\let\BookTitlePage@DefaultTitle\@title
\let\BookTitlePage@DefaultAuthor\@author
\AtBeginDocument{%
  \ifx\@title\BookTitlePage@DefaultTitle\else
    \ifx\BookTitlePage@Edition\@empty
      \hypersetup{pdftitle={\@title}}%
    \else
      \hypersetup{pdftitle={\@title（\BookTitlePage@Edition）}}%
    \fi
  \fi
  \ifx\@author\BookTitlePage@DefaultAuthor\else\hypersetup{pdfauthor={\@author}}\fi}
\makeatother

% 仓库地址：设置后，底部信息区首行为 GitHub 标志加地址，整行链接到仓库；提交哈希链接到 GitHub 上的
% 该提交“<仓库地址>/commit/<短哈希>”，外观不变（hidelinks）。地址不带末尾的 /，按原样读取，
% ~ _ % 等字符不需转义；不写或参数为空时没有地址行，哈希也不加链接。
\ExplSyntaxOn
\str_new:N \g__book_titlepage_repository_str
\NewDocumentCommand \RepositoryURL { v }
  { \str_gset:Nn \g__book_titlepage_repository_str {#1} }
\ExplSyntaxOff

% 提交信息文件由 scripts/commit-info.sh 经根目录 run.sh 写入编译输出目录，
% 四行依次为日期、短哈希、提交说明、clean 或 dirty。按原样逐行读取，
% 不经 TeX 转义；文件不存在或缺少日期、哈希时不显示提交信息。
\ExplSyntaxOn
\ior_new:N \g__book_titlepage_ior
\str_new:N \g__book_titlepage_date_str
\str_new:N \g__book_titlepage_hash_str
\str_new:N \g__book_titlepage_subject_str
\str_new:N \g__book_titlepage_state_str
\bool_new:N \g__book_titlepage_commit_bool
\cs_new_protected:Npn \__book_titlepage_get:N #1
  {
    \ior_if_eof:NTF \g__book_titlepage_ior
      { \str_gclear:N #1 }
      {
        \ior_str_get:NN \g__book_titlepage_ior \l_tmpa_str
        \str_gset_eq:NN #1 \l_tmpa_str
      }
  }
\NewDocumentCommand \TitlePageReadCommit { m }
  {
    \bool_gset_false:N \g__book_titlepage_commit_bool
    \file_if_exist:nT {#1}
      {
        \ior_open:Nn \g__book_titlepage_ior {#1}
        \__book_titlepage_get:N \g__book_titlepage_date_str
        \__book_titlepage_get:N \g__book_titlepage_hash_str
        \__book_titlepage_get:N \g__book_titlepage_subject_str
        \__book_titlepage_get:N \g__book_titlepage_state_str
        \ior_close:N \g__book_titlepage_ior
        \bool_lazy_or:nnF
          { \str_if_empty_p:N \g__book_titlepage_date_str }
          { \str_if_empty_p:N \g__book_titlepage_hash_str }
          { \bool_gset_true:N \g__book_titlepage_commit_bool }
      }
  }
% 默认读取本次编译的 <jobname>.commit；样张可在 \maketitle 前改读其他文件。
\AtBeginDocument { \TitlePageReadCommit { \c_sys_jobname_str .commit } }
% 有提交信息或设了仓库地址时，排底部信息区。
\NewDocumentCommand \BookTitlePageIfStampTF { +m +m }
  {
    \bool_lazy_or:nnTF
      { \g__book_titlepage_commit_bool }
      { ! \str_if_empty_p:N \g__book_titlepage_repository_str }
      {#1} {#2}
  }
% 仓库地址行（用户确认）：Nerd Fonts 的 GitHub 标志 oct-mark_github 加去掉开头 https:// 的地址，
% 与日期同为等宽次要灰；整行链接到仓库，标志与地址作为整体居中。标志取自等宽字体，字形宽 1 em，
% 步进却只有 0.6 em，所以放进 1 em 宽的盒子，再空 0.45 em 接地址。
\str_new:N \l__book_titlepage_repository_text_str
\cs_new_protected:Npn \__book_titlepage_repository_line:
  {
    \str_set_eq:NN \l__book_titlepage_repository_text_str \g__book_titlepage_repository_str
    \str_if_eq:eeT
      { \str_range:Nnn \l__book_titlepage_repository_text_str { 1 } { 8 } } { https:// }
      {
        \str_set:Ne \l__book_titlepage_repository_text_str
          { \str_range:Nnn \l__book_titlepage_repository_text_str { 9 } { -1 } }
      }
    \exp_args:Ne \href { \g__book_titlepage_repository_str }
      {
        \ttfamily
        \makebox [ 1em ] [ l ] { \symbol { "F408 } }
        \hspace { 0.45em }
        \str_use:N \l__book_titlepage_repository_text_str
      }
  }
% 底部信息区：设了仓库地址时首行为地址行；有提交信息时其下为眉题式三行：眉题“LAST COMMIT”；
% 日期与哈希，dirty 时附圆角细框标签；提交说明。不用分隔符串联，依靠换行、留白及字体差异区分标签与数值。
% 只有哈希本身链接到提交；dirty 时同样链接到所基于的提交。
% 间距（用户确认）：行高 11 pt；地址行到眉题的基线间距为 1.5 倍行高，眉题与它标注的两行按行高排成一组，
% 组间空白约为组内的 2 倍。dirty 时标签细框高出数字约 1 pt、低于基线约 2 pt，日期行上下各补足这么多，
% 使细框上下的空白与没有标签时接近。
\bool_new:N \l__book_titlepage_dirty_bool
\NewDocumentCommand \BookTitlePageStampText { }
  {
    \str_if_empty:NF \g__book_titlepage_repository_str
      {
        \__book_titlepage_repository_line:
        \bool_if:NT \g__book_titlepage_commit_bool { \\ [ 5.5pt ] }
      }
    \bool_if:NT \g__book_titlepage_commit_bool
      {
        \str_if_eq:VnTF \g__book_titlepage_state_str { dirty }
          { \bool_set_true:N \l__book_titlepage_dirty_bool }
          { \bool_set_false:N \l__book_titlepage_dirty_bool }
        { \BookTitlePageEyebrow \fontsize { 6.5 } { 8 } \selectfont LAST ~ COMMIT }
        \bool_if:NTF \l__book_titlepage_dirty_bool { \\ [ 1pt ] } { \\ }
        { \ttfamily \str_use:N \g__book_titlepage_date_str }
        \hspace { 0.8em }
        { \ttfamily \color { BookTitlePageAccent }
          \str_if_empty:NTF \g__book_titlepage_repository_str
            { \str_use:N \g__book_titlepage_hash_str }
            {
              \exp_args:Ne \href
                { \g__book_titlepage_repository_str / commit / \g__book_titlepage_hash_str }
                { \str_use:N \g__book_titlepage_hash_str }
            }
        }
        \bool_if:NT \l__book_titlepage_dirty_bool
          { \hspace { 0.6em } \usebox \BookTitlePageDirtyBox }
        \str_if_empty:NF \g__book_titlepage_subject_str
          {
            \bool_if:NTF \l__book_titlepage_dirty_bool { \\ [ 2pt ] } { \\ }
            \BookTitlePagePlainLatin \str_use:N \g__book_titlepage_subject_str
          }
      }
  }
\ExplSyntaxOff
% 未提交修改标签：0.4 pt 浅灰圆角细框，与代码卡片同一视觉语言。
% 在扉页的 tikzpicture 之外预先排入盒子，避免嵌套 TikZ。
\newsavebox{\BookTitlePageDirtyBox}
\newcommand*{\BookTitlePageMakeDirtyBox}{%
  \sbox{\BookTitlePageDirtyBox}{\tikz[baseline=(tag.base)]
    \node[draw=BookTitlePageHairline,line width=0.4pt,rounded corners=1.1mm,
      inner xsep=1.3mm,inner ysep=0.45mm,text=BookTitlePageInk,
      font=\sffamily\fontsize{6.5}{8}\selectfont] (tag) {含未提交修改};}}

% 版面以纸张中心为原点：标题区在上，配图居中偏下，底部信息区贴近底边；双栏正文不影响绝对定位的扉页。
% 版本标记（用户确认，2026-10-15）：全角括号、苹方常规体 22 pt（书名字号的一半）、
% 次要灰，与书名同基线，悬挂在书名右侧、不参与居中；书名仍以页面中轴居中，
% 与英文书名、作者对齐。
\makeatletter
\renewcommand{\maketitle}{%
  \clearpage
  \thispagestyle{empty}%
  \BookTitlePageMakeDirtyBox
  \null
  \begin{tikzpicture}[remember picture,overlay]
    \begin{scope}[shift={(current page.center)}]
      \node[anchor=base,inner sep=0pt,text=BookTitlePageInk,
        font=\sffamily\bfseries\fontsize{44}{52}\selectfont] (title) at (0,\BookTitlePageTitleY) {\@title};
      \ifx\BookTitlePage@Edition\@empty\else
        \node[anchor=base west,inner sep=0pt,text=BookTitlePageSecondary,
          font=\sffamily\mdseries\fontsize{22}{26}\selectfont] at (title.base east)
          {（\BookTitlePage@Edition）};
      \fi
      \ifx\BookTitlePage@EnglishTitle\@empty\else
        \node[anchor=base,inner sep=0pt,text=BookTitlePageSecondary,
          font=\sffamily\fontsize{15}{18}\selectfont] at (0,\BookTitlePageEnglishY)
          {\BookTitlePage@EnglishTitle};
      \fi
      \node[anchor=base,inner sep=0pt,text=BookTitlePageInk,
        font=\BookTitlePageMediumLatin\BookTitlePageMediumCJK\fontsize{12}{15}\selectfont]
        at (0,\BookTitlePageAuthorY) {\@author};
      \ifx\BookTitlePage@Image\@empty\else
        \node[inner sep=0pt] at (0,\BookTitlePageImageY) {\includegraphics[
          width=\BookTitlePageImageWidth,height=\BookTitlePageImageHeight,
          keepaspectratio]{\BookTitlePage@Image}};
      \fi
      % 底边对齐：整块从底边向上排，地址行在提交信息之上，提交说明换行也向上增长，不进入下边距。
      % flush center 使用 fil 胶居中；align=center 的有限伸缩会拉开 xeCJK 字距。
      \BookTitlePageIfStampTF{%
        \node[anchor=south,inner sep=0pt,text width=150mm,align=flush center,
          text=BookTitlePageSecondary,font=\sffamily\fontsize{8}{11}\selectfont]
          at (0,\BookTitlePageStampY) {\BookTitlePageStampText};
      }{}%
    \end{scope}
  \end{tikzpicture}%
  \clearpage}
\makeatother

\TitlePageSetup{title=57mm, english=43.5mm, author=32mm, image=-17mm, stamp=-88mm,
  image-width=110mm, image-height=70mm}


% 行内代码保持 C07 的 \texttt / \verb 语法，只统一字号与字体族。
\let\PortraitOriginalTexttt\texttt
\RenewDocumentCommand{\texttt}{m}{{\fontsize{8.1bp}{14bp}\selectfont\PortraitOriginalTexttt{#1}}}
\makeatletter
\def\verbatim@font{\normalfont\ttfamily\fontsize{8.1bp}{14bp}\selectfont}
\makeatother

% xcode 为试排候选，可按语言覆盖；源代码与缩进不变。
\newcommand\CodeFontSize{\fontsize{8.1bp}{12.5bp}\selectfont}
\definecolor{CodeNumber}{HTML}{999999}
\setminted{style=xcode,fontsize=\CodeFontSize,baselinestretch=1,
  linenos,firstnumber=1,numbersep=5bp,xleftmargin=22bp,
  tabsize=4,autogobble=false,mathescape=false,texcomments=false,
  showspaces=false,showtabs=false,breaklines=true,breakanywhere=true,
  breaksymbolleft={\textcolor{CodeNumber}{\ensuremath{\hookrightarrow}}},
  breaksymbolright={},breakindent=0pt}
\renewcommand\theFancyVerbLine{\textcolor{CodeNumber}{\ttfamily\fontsize{7bp}{9bp}\selectfont\arabic{FancyVerbLine}}}

% 与 HTML <mark> 相同的纯黄底色；零内边距，不改变短语尺寸。
% \colorbox 把短语作为不可拆分的行内盒子，局部分组防止间距设置外泄。
\newcommand{\Highlight}[1]{%
  \begingroup
  \setlength{\fboxsep}{0pt}%
  \colorbox[HTML]{FFFF00}{#1}%
  \endgroup
}

% 代码标题栏固定外观；代码字体、配色、行号与折行由 \setminted[语言] 控制。
\definecolor{InputCodeTitle}{HTML}{F2F2F2}
\ExplSyntaxOn
\str_new:N \l__inputcode_area_str
\str_new:N \l__inputcode_base_str
\str_new:N \l__inputcode_ext_str
\str_new:N \l__inputcode_title_str

\cs_new_protected:Npn \__inputcode_filename:n #1
  {
    \file_parse_full_name:nNNN {#1}
      \l__inputcode_area_str \l__inputcode_base_str \l__inputcode_ext_str
    \str_set:Nx \l__inputcode_title_str
      {\str_use:N \l__inputcode_base_str \str_use:N \l__inputcode_ext_str}
  }

\cs_new_protected:Npn \__inputcode_print_title:n #1
  {
    \str_map_inline:nn {#1}
      {
        \str_if_eq:nnTF {##1} { ~ }
          {\hspace{\fontdimen2\font}}
          {##1}
      }
  }

\cs_new_protected:Npn \__inputcode_box:nnn #1#2#3
  {
    \par
    \begin{tcolorbox}[
      enhanced,
      breakable,
      title={{\xeCJKsetup{CJKecglue={}}\__inputcode_print_title:n {#1}}},
      colback=white,
      colbacktitle=InputCodeTitle,
      colframe=black,
      coltitle=black,
      boxrule=0.4pt,
      titlerule=0.4pt,
      arc=0pt,
      left=1mm,right=1mm,top=1mm,bottom=1mm,
      lefttitle=1mm,righttitle=1mm,toptitle=0.5mm,bottomtitle=0.5mm,
      fonttitle={\normalfont\ttfamily\mdseries\upshape\fontsize{9bp}{14bp}\selectfont}
    ]
    \xeCJKsetup{CJKecglue={}}
    \inputminted[mathescape=false,texcomments=false,autogobble=false]{#2}{#3}
    \end{tcolorbox}
  }

\NewDocumentCommand{\InputCode}{o m m}
  {
    \IfNoValueTF{#1}
      {
        \__inputcode_filename:n {#3}
        \exp_args:NV \__inputcode_box:nnn \l__inputcode_title_str {#2} {#3}
      }
      {\__inputcode_box:nnn {#1} {#2} {#3}}
  }
\ExplSyntaxOff

% 可跨页的引用块：浅灰背景、左侧深灰竖线，无其余边框。
\definecolor{QuoteBlockBackground}{HTML}{F7F7F7}
\definecolor{QuoteBlockRule}{HTML}{777777}
\newtcolorbox{quoteblock}{
  enhanced,
  breakable,
  colback=QuoteBlockBackground,
  boxrule=0pt,
  borderline west={2pt}{0pt}{QuoteBlockRule},
  arc=0pt,
  boxsep=0pt,
  left=2mm,right=2mm,top=2mm,bottom=2mm
}

% Markdown 水平分隔线：按当前容器的 \linewidth 排版，不改变分页。
\definecolor{SeparatorRule}{HTML}{A0A0A0}
\NewDocumentCommand{\Separator}{}{%
  \par\addvspace{\baselineskip}%
  \noindent{\color{SeparatorRule}\rule{\linewidth}{0.4pt}}\par
  \addvspace{\baselineskip}%
}


\title{算法竞赛小记}
\author{Calculatelove}
\date{}
\EnglishTitle{Competitive Programming Library}
\TitlePageImage{assets/cl-logo-c4.png}
\EditionLabel{代码版}

\begin{document}

\maketitle

\tableofcontents

% 来源：contents/0x10 基础/0x11 基础 Misc.tex:1
\section{0x11\ 基础\ Misc}

% 来源：contents/0x10 基础/0x11 基础 Misc.tex:11
\subsection{常用公式}

% 来源：contents/0x10 基础/0x11 基础 Misc.tex:13
\subsubsection{日期公式}

% 来源：contents/0x10 基础/0x11 基础 Misc.tex:15
\InputCode{cpp}{contents/0x10 基础/code/0x11 日期公式.cpp}

% 来源：contents/0x10 基础/0x11 基础 Misc.tex:17
\subsubsection{xorshift}

% 来源：contents/0x10 基础/0x11 基础 Misc.tex:19
\InputCode{cpp}{contents/0x10 基础/code/0x11 xorshift.cpp}

% 来源：contents/0x10 基础/0x11 基础 Misc.tex:21
\subsubsection{xorsum}

% 来源：contents/0x10 基础/0x11 基础 Misc.tex:31
\InputCode{cpp}{contents/0x10 基础/code/0x11 xorsum.cpp}

% 来源：contents/0x10 基础/0x11 基础 Misc.tex:42
\subsection{常用配置}

% 来源：contents/0x10 基础/0x11 基础 Misc.tex:44
\subsubsection{makefile}

% 来源：contents/0x10 基础/0x11 基础 Misc.tex:46
\begin{minted}{make}
CXXFLAGS = -std=c++23 -O2 -Wall -fsanitize=undefined

# -D_GLIBCXX_DEBUG : STL debug mode
# -fsanitize=address : 内存错误检查
# -fsanitize=undefined : UB 检查
\end{minted}

% 来源：contents/0x10 基础/0x11 基础 Misc.tex:54
\subsubsection{pragma}

% 来源：contents/0x10 基础/0x11 基础 Misc.tex:56
\begin{minted}{cpp}
#pragma GCC optimize("O3,unroll-loops")
#pragma GCC target("avx2")
\end{minted}

% 来源：contents/0x10 基础/0x11 基础 Misc.tex:61
\subsubsection{ulimit}

% 来源：contents/0x10 基础/0x11 基础 Misc.tex:63
\begin{minted}{bash}
ulimit -s 1048576
\end{minted}

% 来源：contents/0x10 基础/0x11 基础 Misc.tex:67
\subsubsection{Linux 对拍脚本}

% 来源：contents/0x10 基础/0x11 基础 Misc.tex:71
\begin{minted}{bash}
while true; do
    ./gen >0.in
    ./obj <0.in >1.out
    ./std <0.in >2.out
    let cnt=cnt+1; echo "Test $cnt:"
    if diff 1.out 2.out; then
        echo ac
    else
        echo wa
        break;
    fi
done
\end{minted}

% 来源：contents/0x10 基础/0x11 基础 Misc.tex:86
\subsubsection{Windows 对拍程序}

% 来源：contents/0x10 基础/0x11 基础 Misc.tex:88
\begin{minted}{cpp}
#include <bits/stdc++.h>
int main() {
    int T = 100;
    for (int i = 1; i <= T; i ++) {
        system("gen.exe >0.in");
        system("obj.exe <0.in >1.out");
        system("std.exe <0.in >2.out");
        std::cout << "Test " << i << ":\n";
        if (system("fc 1.out 2.out")) {
            std::cout << "wa\n";
            break;
        } else {
            std::cout << "ac\n";
        }
    }
}
\end{minted}

% 来源：contents/0x10 基础/0x11 基础 Misc.tex:107
\subsection{随机数据生成}

% 来源：contents/0x10 基础/0x11 基础 Misc.tex:109
\subsubsection{生成区间 $[l,r]$ 中的随机整数}

% 来源：contents/0x10 基础/0x11 基础 Misc.tex:111
\InputCode{cpp}{contents/0x10 基础/code/0x11 生成区间 [l,r] 中的随机整数.cpp}

% 来源：contents/0x10 基础/0x11 基础 Misc.tex:113
\subsubsection{生成区间 $[l,r)$ 中的随机浮点数}

% 来源：contents/0x10 基础/0x11 基础 Misc.tex:115
\InputCode{cpp}{contents/0x10 基础/code/0x11 生成区间 [l,r) 中的随机浮点数.cpp}

% 来源：contents/0x10 基础/0x11 基础 Misc.tex:117
\subsubsection{生成互不相同的随机整数}

% 来源：contents/0x10 基础/0x11 基础 Misc.tex:119
\InputCode{cpp}{contents/0x10 基础/code/0x11 生成互不相同的随机整数.cpp}

% 来源：contents/0x10 基础/0x11 基础 Misc.tex:121
\subsubsection{生成随机树}

% 来源：contents/0x10 基础/0x11 基础 Misc.tex:127
\InputCode{cpp}{contents/0x10 基础/code/0x11 生成随机树.cpp}

% 来源：contents/0x10 基础/0x11 基础 Misc.tex:129
\subsection{常用库函数重载}

% 来源：contents/0x10 基础/0x11 基础 Misc.tex:131
\subsubsection{i128}

% 来源：contents/0x10 基础/0x11 基础 Misc.tex:135
\InputCode{cpp}{contents/0x10 基础/code/0x11 i128.cpp}

% 来源：contents/0x10 基础/0x11 基础 Misc.tex:137
\subsubsection{整数除法精确取整}

% 来源：contents/0x10 基础/0x11 基础 Misc.tex:139
\InputCode{cpp}{contents/0x10 基础/code/0x11 整数除法精确取整.cpp}

% 来源：contents/0x10 基础/0x11 基础 Misc.tex:141
\subsubsection{取最值}

% 来源：contents/0x10 基础/0x11 基础 Misc.tex:143
\InputCode{cpp}{contents/0x10 基础/code/0x11 取最值.cpp}

% 来源：contents/0x10 基础/0x11 基础 Misc.tex:145
\subsubsection{取模}

% 来源：contents/0x10 基础/0x11 基础 Misc.tex:147
\InputCode{cpp}{contents/0x10 基础/code/0x11 取模.cpp}

% 来源：contents/0x10 基础/0x11 基础 Misc.tex:149
\subsubsection{精确开根号}

% 来源：contents/0x10 基础/0x11 基础 Misc.tex:151
\InputCode{cpp}{contents/0x10 基础/code/0x11 精确开根号.cpp}

% 来源：contents/0x10 基础/0x11 基础 Misc.tex:153
\subsubsection{精确取对数}

% 来源：contents/0x10 基础/0x11 基础 Misc.tex:155
\InputCode{cpp}{contents/0x10 基础/code/0x11 精确取对数.cpp}

% 来源：contents/0x10 基础/0x11 基础 Misc.tex:157
\subsubsection{快读}

% 来源：contents/0x10 基础/0x11 基础 Misc.tex:159
\InputCode{cpp}{contents/0x10 基础/code/0x11 快读.cpp}

% 来源：contents/0x10 基础/0x11 基础 Misc.tex:161
\subsubsection{超级快读}

% 来源：contents/0x10 基础/0x11 基础 Misc.tex:163
\InputCode{cpp}{contents/0x10 基础/code/0x11 超级快读.cpp}

% 来源：contents/0x10 基础/0x11 基础 Misc.tex:165
\subsection{高精度}

% 来源：contents/0x10 基础/0x11 基础 Misc.tex:167
\InputCode{cpp}{contents/0x10 基础/code/0x11 高精度.cpp}

% 来源：contents/0x10 基础/0x11 基础 Misc.tex:169
\subsection{分数类}

% 来源：contents/0x10 基础/0x11 基础 Misc.tex:171
\InputCode{cpp}{contents/0x10 基础/code/0x11 分数类.cpp}

% 来源：contents/0x10 基础/0x12 基础 算法.tex:1
\section{0x12\ 基础\ 算法}

% 来源：contents/0x10 基础/0x12 基础 算法.tex:47
\subsection{快速乘}

% 来源：contents/0x10 基础/0x12 基础 算法.tex:53
\subsubsection{正道：i128}

% 来源：contents/0x10 基础/0x12 基础 算法.tex:55
\InputCode{cpp}{contents/0x10 基础/code/0x12 快速乘（i128）.cpp}

% 来源：contents/0x10 基础/0x12 基础 算法.tex:57
\subsubsection{魔道：long double}

% 来源：contents/0x10 基础/0x12 基础 算法.tex:63
\InputCode{cpp}{contents/0x10 基础/code/0x12 快速乘（long double）.cpp}

% 来源：contents/0x10 基础/0x12 基础 算法.tex:65
\subsection{光速幂}

% 来源：contents/0x10 基础/0x12 基础 算法.tex:81
\InputCode{cpp}{contents/0x10 基础/code/0x12 光速幂.cpp}

% 来源：contents/0x20 字符串/0x21 字符串 I.tex:1
\section{0x21\ 字符串\ I}

% 来源：contents/0x20 字符串/0x21 字符串 I.tex:94
\subsection{kmp}

% 来源：contents/0x20 字符串/0x21 字符串 I.tex:108
\InputCode{cpp}{contents/0x20 字符串/code/0x21 kmp.cpp}

% 来源：contents/0x20 字符串/0x21 字符串 I.tex:166
\subsection{Z Algorithm}

% 来源：contents/0x20 字符串/0x21 字符串 I.tex:188
\InputCode{cpp}{contents/0x20 字符串/code/0x21 Z-Algorithm.cpp}

% 来源：contents/0x20 字符串/0x21 字符串 I.tex:190
\subsection{manacher}

% 来源：contents/0x20 字符串/0x21 字符串 I.tex:218
\InputCode{cpp}{contents/0x20 字符串/code/0x21 manacher.cpp}

% 来源：contents/0x20 字符串/0x21 字符串 I.tex:220
\InputCode{cpp}{contents/0x20 字符串/code/0x21 manacher（奇偶分开考虑）.cpp}

% 来源：contents/0x20 字符串/0x21 字符串 I.tex:230
\subsection{Lyndon 分解}

% 来源：contents/0x20 字符串/0x21 字符串 I.tex:270
\InputCode{cpp}{contents/0x20 字符串/code/0x21 duval.cpp}

% 来源：contents/0x20 字符串/0x21 字符串 I.tex:272
\subsection{最小表示法}

% 来源：contents/0x20 字符串/0x21 字符串 I.tex:290
\InputCode{cpp}{contents/0x20 字符串/code/0x21 最小表示法.cpp}

% 来源：contents/0x20 字符串/0x22 字符串 II.tex:1
\section{0x22\ 字符串\ II}

% 来源：contents/0x20 字符串/0x22 字符串 II.tex:4
\subsection{Trie}

% 来源：contents/0x20 字符串/0x22 字符串 II.tex:26
\InputCode{cpp}{contents/0x20 字符串/code/0x22 01 trie.cpp}

% 来源：contents/0x20 字符串/0x22 字符串 II.tex:28
\InputCode{cpp}{contents/0x20 字符串/code/0x22 可持久化 01 trie.cpp}

% 来源：contents/0x20 字符串/0x22 字符串 II.tex:71
\InputCode{cpp}{contents/0x20 字符串/code/0x22 01 trie 维护集合异或和.cpp}

% 来源：contents/0x20 字符串/0x22 字符串 II.tex:73
\subsection{AC 自动机}

% 来源：contents/0x20 字符串/0x22 字符串 II.tex:91
\InputCode{cpp}{contents/0x20 字符串/code/0x22 AC 自动机.cpp}

% 来源：contents/0x20 字符串/0x22 字符串 II.tex:95
\InputCode{cpp}{contents/0x20 字符串/code/0x22 AC 自动机（大字符集）.cpp}

% 来源：contents/0x20 字符串/0x22 字符串 II.tex:107
\subsection{SA}

% 来源：contents/0x20 字符串/0x22 字符串 II.tex:137
\InputCode{cpp}{contents/0x20 字符串/code/0x22 SA.cpp}

% 来源：contents/0x20 字符串/0x22 字符串 II.tex:139
\subsection{SAM}

% 来源：contents/0x20 字符串/0x22 字符串 II.tex:177
\InputCode{cpp}{contents/0x20 字符串/code/0x22 SAM.cpp}

% 来源：contents/0x20 字符串/0x22 字符串 II.tex:189
\begin{minted}{cpp}
int p = 1, len = 0; // 当前匹配到的状态与长度
for (int i = 1; i <= m; i ++) {
    int c = T[i] - 'a';
    while (p > 1 && !SAM::t[p].trans[c]) {
        p = SAM::t[p].link, len = SAM::t[p].maxl;
    }
    if (SAM::t[p].trans[c]) {
        p = SAM::t[p].trans[c], len ++;
    }
    // 此时的 p 与 len，就是以 i 为结尾时匹配到的最长状态与长度
}
\end{minted}

% 来源：contents/0x20 字符串/0x22 字符串 II.tex:222
\InputCode{cpp}{contents/0x20 字符串/code/0x22 SAM 前端操作.cpp}

% 来源：contents/0x20 字符串/0x22 字符串 II.tex:234
\InputCode{cpp}{contents/0x20 字符串/code/0x22 广义 SAM（离线）.cpp}

% 来源：contents/0x20 字符串/0x22 字符串 II.tex:236
\InputCode{cpp}{contents/0x20 字符串/code/0x22 广义 SAM（在线）.cpp}

% 来源：contents/0x20 字符串/0x22 字符串 II.tex:238
\subsection{PAM}

% 来源：contents/0x20 字符串/0x22 字符串 II.tex:254
\InputCode{cpp}{contents/0x20 字符串/code/0x22 PAM.cpp}

% 来源：contents/0x20 字符串/0x22 字符串 II.tex:258
\InputCode{cpp}{contents/0x20 字符串/code/0x22 PAM half.cpp}

% 来源：contents/0x20 字符串/0x22 字符串 II.tex:260
\subsubsection{Palindrome Series}

% 来源：contents/0x20 字符串/0x22 字符串 II.tex:307
\InputCode{cpp}{contents/0x20 字符串/code/0x22 PAM series-link.cpp}

% 来源：contents/0x30 数学/0x31 数学 多项式.tex:1
\section{0x31\ 数学\ 多项式}

% 来源：contents/0x30 数学/0x31 数学 多项式.tex:4
\subsection{拉格朗日插值}

% 来源：contents/0x30 数学/0x31 数学 多项式.tex:25
\InputCode{cpp}{contents/0x30 数学/code/0x31 拉格朗日插值.cpp}

% 来源：contents/0x30 数学/0x31 数学 多项式.tex:43
\InputCode{cpp}{contents/0x30 数学/code/0x31 拉格朗日插值求系数.cpp}

% 来源：contents/0x30 数学/0x31 数学 多项式.tex:70
\InputCode{cpp}{contents/0x30 数学/code/0x31 连续点值的拉格朗日插值.cpp}

% 来源：contents/0x30 数学/0x31 数学 多项式.tex:72
\subsection{FFT}

% 来源：contents/0x30 数学/0x31 数学 多项式.tex:270
\subsubsection{FFT 多项式乘法}

% 来源：contents/0x30 数学/0x31 数学 多项式.tex:282
\InputCode{cpp}{contents/0x30 数学/code/0x31 FFT.cpp}

% 来源：contents/0x30 数学/0x31 数学 多项式.tex:284
\subsection{NTT}

% 来源：contents/0x30 数学/0x31 数学 多项式.tex:336
\InputCode{cpp}{contents/0x30 数学/code/0x31 NTT.cpp}

% 来源：contents/0x30 数学/0x31 数学 多项式.tex:338
\subsection{MTT}

% 来源：contents/0x30 数学/0x31 数学 多项式.tex:344
\subsubsection{三模 NTT}

% 来源：contents/0x30 数学/0x31 数学 多项式.tex:358
\InputCode{cpp}{contents/0x30 数学/code/0x31 MTT（三模 NTT）.cpp}

% 来源：contents/0x30 数学/0x31 数学 多项式.tex:811
\subsection{多项式多点求值}

% 来源：contents/0x30 数学/0x31 数学 多项式.tex:835
\InputCode{cpp}{contents/0x30 数学/code/0x31 多项式多点求值（多项式取模）.cpp}

% 来源：contents/0x30 数学/0x31 数学 多项式.tex:927
\subsection{多项式全家桶⭐}

% 来源：contents/0x30 数学/0x31 数学 多项式.tex:942
\InputCode{cpp}{contents/0x30 数学/code/0x31 多项式全家桶.cpp}

% 来源：contents/0x30 数学/0x31 数学 多项式.tex:944
\subsection{Bostan-Mori}

% 来源：contents/0x30 数学/0x31 数学 多项式.tex:946
\subsubsection{分式远处系数}

% 来源：contents/0x30 数学/0x31 数学 多项式.tex:972
\InputCode{cpp}{contents/0x30 数学/code/0x31 分式远处系数.cpp}

% 来源：contents/0x30 数学/0x31 数学 多项式.tex:978
\subsection{常系数齐次线性递推}

% 来源：contents/0x30 数学/0x31 数学 多项式.tex:982
\subsubsection{Fiduccia 算法}

% 来源：contents/0x30 数学/0x31 数学 多项式.tex:998
\InputCode{cpp}{contents/0x30 数学/code/0x31 常系数齐次线性递推（Fiduccia）.cpp}

% 来源：contents/0x30 数学/0x31 数学 多项式.tex:1000
\subsubsection{递推式转分式}

% 来源：contents/0x30 数学/0x31 数学 多项式.tex:1023
\InputCode{cpp}{contents/0x30 数学/code/0x31 常系数齐次线性递推（Bostan-Mori）.cpp}

% 来源：contents/0x30 数学/0x31 数学 多项式.tex:1025
\subsection{整式递推}

% 来源：contents/0x30 数学/0x31 数学 多项式.tex:1029
\subsubsection{应用：快速阶乘算法}

% 来源：contents/0x30 数学/0x31 数学 多项式.tex:1065
\InputCode{cpp}{contents/0x30 数学/code/0x31 快速阶乘算法.cpp}

% 来源：contents/0x30 数学/0x32 数学 集合幂级数.tex:1
\section{0x32\ 数学\ 集合幂级数}

% 来源：contents/0x30 数学/0x32 数学 集合幂级数.tex:4
\subsection{FWT}

% 来源：contents/0x30 数学/0x32 数学 集合幂级数.tex:106
\subsubsection{位运算卷积}

% 来源：contents/0x30 数学/0x32 数学 集合幂级数.tex:146
\InputCode{cpp}{contents/0x30 数学/code/0x32 FWT.cpp}

% 来源：contents/0x30 数学/0x32 数学 集合幂级数.tex:148
\subsection{子集 dp}

% 来源：contents/0x30 数学/0x32 数学 集合幂级数.tex:158
\begin{minted}{cpp}
// 注意：要先枚举维度 i，再枚举状态 S
for (int i = 0; i < m; i ++) {
    for (int S = 0; S < (1 << m); S ++) {
        if (S >> i & 1) {
            f[S] += f[S ^ (1 << i)];
        }
    }
}
\end{minted}

% 来源：contents/0x30 数学/0x32 数学 集合幂级数.tex:175
\begin{minted}{cpp}
for (int s = 1; s <= m; s ++) {
    for (int S = 0; S < (1 << m); S ++) {
        tmp[S] = f[S];
    }
    for (int i = 0; i < m; i ++) {
        for (int S = 0; S < (1 << m); S ++) {
            if (S >> i & 1) {
                tmp[S] += tmp[S ^ (1 << i)];
            }
        }
    }
    
    for (int S = 0; S < (1 << m); S ++) {
        if (__builtin_popcount(S) == s) {
            // 使用 tmp[S] 更新 f[S]
        }
    }
}
\end{minted}

% 来源：contents/0x30 数学/0x32 数学 集合幂级数.tex:196
\subsection{子集卷积}

% 来源：contents/0x30 数学/0x32 数学 集合幂级数.tex:221
\InputCode{cpp}{contents/0x30 数学/code/0x32 子集卷积.cpp}

% 来源：contents/0x30 数学/0x32 数学 集合幂级数.tex:233
\InputCode{cpp}{contents/0x30 数学/code/0x32 半在线子集卷积.cpp}

% 来源：contents/0x30 数学/0x33 数学 组合计数.tex:1
\section{0x33\ 数学\ 组合计数}

% 来源：contents/0x30 数学/0x33 数学 组合计数.tex:139
\subsection{组合数常见求法}

% 来源：contents/0x30 数学/0x33 数学 组合计数.tex:141
\subsubsection{求法 1：递推}

% 来源：contents/0x30 数学/0x33 数学 组合计数.tex:151
\begin{minted}{cpp}
for (int i = 0; i <= n; i ++) {
    C[i][0] = C[i][i] = 1;
    for (int j = 1; j < i; j ++) {
        C[i][j] = (C[i - 1][j] + C[i - 1][j - 1]) % mod;
    }
}
\end{minted}

% 来源：contents/0x30 数学/0x33 数学 组合计数.tex:160
\subsubsection{求法 2：预处理阶乘与阶乘逆元}

% 来源：contents/0x30 数学/0x33 数学 组合计数.tex:180
\InputCode{cpp}{contents/0x30 数学/code/0x33 组合数（预处理阶乘与阶乘逆元）.cpp}

% 来源：contents/0x30 数学/0x33 数学 组合计数.tex:634
\subsection{图论计数}

% 来源：contents/0x30 数学/0x33 数学 组合计数.tex:636
\subsubsection{有标号无向连通图计数}

% 来源：contents/0x30 数学/0x33 数学 组合计数.tex:661
\paragraph{算法二：多项式 ln}

% 来源：contents/0x30 数学/0x33 数学 组合计数.tex:669
\InputCode{cpp}{contents/0x30 数学/code/0x33 有标号无向连通图计数.cpp}

% 来源：contents/0x30 数学/0x33 数学 组合计数.tex:671
\subsubsection{有标号无向欧拉图计数}

% 来源：contents/0x30 数学/0x33 数学 组合计数.tex:681
\InputCode{cpp}{contents/0x30 数学/code/0x33 有标号无向欧拉图计数.cpp}

% 来源：contents/0x30 数学/0x33 数学 组合计数.tex:683
\subsubsection{有标号无向二分图计数}

% 来源：contents/0x30 数学/0x33 数学 组合计数.tex:715
\InputCode{cpp}{contents/0x30 数学/code/0x33 有标号无向二分图计数.cpp}

% 来源：contents/0x30 数学/0x33 数学 组合计数.tex:717
\subsubsection{有标号 DAG 计数}

% 来源：contents/0x30 数学/0x33 数学 组合计数.tex:740
\InputCode{cpp}{contents/0x30 数学/code/0x33 有标号 DAG 计数.cpp}

% 来源：contents/0x30 数学/0x33 数学 组合计数.tex:742
\subsubsection{有标号弱连通 DAG 计数}

% 来源：contents/0x30 数学/0x33 数学 组合计数.tex:752
\InputCode{cpp}{contents/0x30 数学/code/0x33 有标号弱连通 DAG 计数.cpp}

% 来源：contents/0x30 数学/0x33 数学 组合计数.tex:754
\subsubsection{有标号强连通图计数}

% 来源：contents/0x30 数学/0x33 数学 组合计数.tex:797
\InputCode{cpp}{contents/0x30 数学/code/0x33 有标号强连通图计数.cpp}

% 来源：contents/0x30 数学/0x34 数学 初等数论.tex:1
\section{0x34\ 数学\ 初等数论}

% 来源：contents/0x30 数学/0x34 数学 初等数论.tex:4
\subsection{质数🚩}

% 来源：contents/0x30 数学/0x34 数学 初等数论.tex:39
\subsubsection{埃氏筛}

% 来源：contents/0x30 数学/0x34 数学 初等数论.tex:63
\InputCode{cpp}{contents/0x30 数学/code/0x34 埃氏筛.cpp}

% 来源：contents/0x30 数学/0x34 数学 初等数论.tex:65
\subsubsection{线性筛}

% 来源：contents/0x30 数学/0x34 数学 初等数论.tex:75
\InputCode{cpp}{contents/0x30 数学/code/0x34 线性筛.cpp}

% 来源：contents/0x30 数学/0x34 数学 初等数论.tex:77
\subsection{狄利克雷前缀和}

% 来源：contents/0x30 数学/0x34 数学 初等数论.tex:83
\begin{minted}{cpp}
// 注意：要先枚举质数 prime[j]
for (int j = 1; j <= primeCount; j ++) {
    for (int i = 1; i <= n / prime[j]; i ++) {
        f[i * prime[j]] += f[i];
    }
}
\end{minted}

% 来源：contents/0x30 数学/0x34 数学 初等数论.tex:94
\begin{minted}{cpp}
// 注意：要先枚举质数 prime[j]
for (int j = 1; j <= primeCount; j ++) {
    for (int i = n / prime[j]; i >= 1; i --) {
        f[i] += f[i * prime[j]];
    }
}
\end{minted}

% 来源：contents/0x30 数学/0x34 数学 初等数论.tex:105
\begin{minted}{cpp}
// 注意：要先枚举质数 prime[j]
for (int j = 1; j <= primeCount; j ++) {
    for (int i = n / prime[j]; i >= 1; i --) { // 顺序与正变换相反
        f[i * prime[j]] -= f[i];
    }
}
\end{minted}

% 来源：contents/0x30 数学/0x34 数学 初等数论.tex:116
\begin{minted}{cpp}
// 注意：要先枚举质数 prime[j]
for (int j = 1; j <= primeCount; j ++) {
    for (int i = 1; i <= n / prime[j]; i ++) { // 顺序与正变换相反
        f[i] -= f[i * prime[j]];
    }
}
\end{minted}

% 来源：contents/0x30 数学/0x34 数学 初等数论.tex:161
\subsection{Pollard-Rho 质因数分解}

% 来源：contents/0x30 数学/0x34 数学 初等数论.tex:204
\InputCode{cpp}{contents/0x30 数学/code/0x34 素性测试与质因数分解（Miller-Rabin + Pollard-Rho）.cpp}

% 来源：contents/0x30 数学/0x34 数学 初等数论.tex:334
\subsection{最大公约数常见求法}

% 来源：contents/0x30 数学/0x34 数学 初等数论.tex:338
\subsubsection{求法 1：辗转相除法}

% 来源：contents/0x30 数学/0x34 数学 初等数论.tex:356
\begin{minted}{cpp}
// 辗转相除法
int gcd(int a, int b) {
    return b == 0 ? a : gcd(b, a % b);
}
\end{minted}

% 来源：contents/0x30 数学/0x34 数学 初等数论.tex:363
\subsubsection{求法 2：更相减损术}

% 来源：contents/0x30 数学/0x34 数学 初等数论.tex:388
\InputCode{cpp}{contents/0x30 数学/code/0x34 gcd（更相减损术）.cpp}

% 来源：contents/0x30 数学/0x34 数学 初等数论.tex:390
\subsubsection{求法 3：递推}

% 来源：contents/0x30 数学/0x34 数学 初等数论.tex:396
\begin{minted}{cpp}
for (int i = 0; i <= n; i ++) {
    g[i][0] = g[0][i] = i;
    for (int j = 1; j <= i; j ++) {
        g[i][j] = g[j][i] = g[i - j][j];
    }
}
\end{minted}

% 来源：contents/0x30 数学/0x34 数学 初等数论.tex:405
\subsubsection{求法 4：基于值域预处理的快速 gcd}

% 来源：contents/0x30 数学/0x34 数学 初等数论.tex:437
\InputCode{cpp}{contents/0x30 数学/code/0x34 gcd（基于值域预处理的快速 gcd）.cpp}

% 来源：contents/0x30 数学/0x34 数学 初等数论.tex:439
\subsection{类欧几里得算法}

% 来源：contents/0x30 数学/0x34 数学 初等数论.tex:473
\InputCode{cpp}{contents/0x30 数学/code/0x34 类欧几里得算法.cpp}

% 来源：contents/0x30 数学/0x34 数学 初等数论.tex:515
\subsection{数论分块}

% 来源：contents/0x30 数学/0x34 数学 初等数论.tex:519
\subsubsection{下取整数论分块}

% 来源：contents/0x30 数学/0x34 数学 初等数论.tex:548
\begin{minted}{cpp}
// h 表示 x 的上界
for (int x = 1, nx; x <= h; x = nx + 1) {
    nx = h;
    if (x <= n) {
        chmin(nx, n / (n / x));
    }
    // 区间 [x, nx] 内：下取整的值均为 n/x
}
\end{minted}

% 来源：contents/0x30 数学/0x34 数学 初等数论.tex:563
\subsubsection{上取整数论分块}

% 来源：contents/0x30 数学/0x34 数学 初等数论.tex:571
\begin{minted}{cpp}
// h 表示 x 的上界
for (int x = 1, nx; x <= h; x = nx + 1) {
    nx = h;
    if (x < n) {
        chmin(nx, (n - 1) / ((n - 1) / x));
    }
    // 区间 [x, nx] 内：上取整的值均为 (n-1)/x+1
}
\end{minted}

% 来源：contents/0x30 数学/0x34 数学 初等数论.tex:582
\subsubsection{多维数论分块}

% 来源：contents/0x30 数学/0x34 数学 初等数论.tex:592
\begin{minted}{cpp}
// h 表示 x 的上界
for (int x = 1, nx; x <= h; x = nx + 1) {
    nx = h;
    if (x <= n) {
        chmin(nx, n / (n / x));
    }
    if (x <= m) {
        chmin(nx, m / (m / x));
    }
    // 区间 [x, nx] 内：n, m 对 x 下取整的值分别为 n/x, m/x
}
\end{minted}

% 来源：contents/0x30 数学/0x34 数学 初等数论.tex:606
\subsection{积性函数与狄利克雷卷积}

% 来源：contents/0x30 数学/0x34 数学 初等数论.tex:709
\subsubsection{线性筛筛积性函数}

% 来源：contents/0x30 数学/0x34 数学 初等数论.tex:723
\InputCode{cpp}{contents/0x30 数学/code/0x34 线性筛筛积性函数.cpp}

% 来源：contents/0x30 数学/0x34 数学 初等数论.tex:801
\subsection{杜教筛}

% 来源：contents/0x30 数学/0x34 数学 初等数论.tex:852
\InputCode{cpp}{contents/0x30 数学/code/0x34 杜教筛.cpp}

% 来源：contents/0x30 数学/0x34 数学 初等数论.tex:963
\subsection{模意义下逆元}

% 来源：contents/0x30 数学/0x34 数学 初等数论.tex:975
\subsubsection{单个逆元的求解}

% 来源：contents/0x30 数学/0x34 数学 初等数论.tex:981
\InputCode{cpp}{contents/0x30 数学/code/0x34 单个逆元的求解.cpp}

% 来源：contents/0x30 数学/0x34 数学 初等数论.tex:989
\subsubsection{预处理线性逆元}

% 来源：contents/0x30 数学/0x34 数学 初等数论.tex:1003
\InputCode{cpp}{contents/0x30 数学/code/0x34 预处理线性逆元.cpp}

% 来源：contents/0x30 数学/0x34 数学 初等数论.tex:1005
\subsection{线性同余方程}

% 来源：contents/0x30 数学/0x34 数学 初等数论.tex:1047
\InputCode{cpp}{contents/0x30 数学/code/0x34 exgcd（扩展欧几里得算法）.cpp}

% 来源：contents/0x30 数学/0x34 数学 初等数论.tex:1087
\subsection{线性同余方程组}

% 来源：contents/0x30 数学/0x34 数学 初等数论.tex:1117
\InputCode{cpp}{contents/0x30 数学/code/0x34 CRT（中国剩余定理）.cpp}

% 来源：contents/0x30 数学/0x34 数学 初等数论.tex:1138
\InputCode{cpp}{contents/0x30 数学/code/0x34 exCRT（扩展中国剩余定理）.cpp}

% 来源：contents/0x30 数学/0x34 数学 初等数论.tex:1140
\subsection{阶与原根}

% 来源：contents/0x30 数学/0x34 数学 初等数论.tex:1257
\InputCode{cpp}{contents/0x30 数学/code/0x34 最小原根.cpp}

% 来源：contents/0x30 数学/0x34 数学 初等数论.tex:1259
\InputCode{cpp}{contents/0x30 数学/code/0x34 所有原根.cpp}

% 来源：contents/0x30 数学/0x34 数学 初等数论.tex:1276
\subsection{高次同余方程}

% 来源：contents/0x30 数学/0x34 数学 初等数论.tex:1292
\InputCode{cpp}{contents/0x30 数学/code/0x34 BSGS（大步小步算法）.cpp}

% 来源：contents/0x30 数学/0x34 数学 初等数论.tex:1308
\InputCode{cpp}{contents/0x30 数学/code/0x34 exBSGS（扩展大步小步算法）.cpp}

% 来源：contents/0x40 数据结构/0x41 数据结构 序列结构.tex:1
\section{0x41\ 数据结构\ 序列结构}

% 来源：contents/0x40 数据结构/0x41 数据结构 序列结构.tex:4
\subsection{线段树}

% 来源：contents/0x40 数据结构/0x41 数据结构 序列结构.tex:10
\subsubsection{势能线段树}

% 来源：contents/0x40 数据结构/0x41 数据结构 序列结构.tex:43
\InputCode{cpp}{contents/0x40 数据结构/code/0x41 吉老师线段树.cpp}

% 来源：contents/0x40 数据结构/0x42 数据结构 树上结构.tex:1
\section{0x42\ 数据结构\ 树上结构}

% 来源：contents/0x40 数据结构/0x42 数据结构 树上结构.tex:4
\subsection{并查集}

% 来源：contents/0x40 数据结构/0x42 数据结构 树上结构.tex:6
\InputCode{cpp}{contents/0x40 数据结构/code/0x42 可撤销并查集.cpp}

% 来源：contents/0x60 图论/0x65 图论 二分图与网络流.tex:1
\section{0x65\ 图论\ 二分图与网络流}

% 来源：contents/0x60 图论/0x65 图论 二分图与网络流.tex:4
\subsection{最大流}

% 来源：contents/0x60 图论/0x65 图论 二分图与网络流.tex:6
\subsubsection{Dinic}

% 来源：contents/0x60 图论/0x65 图论 二分图与网络流.tex:8
\InputCode{cpp}{contents/0x60 图论/code/0x65 最大流（Dinic）.cpp}

\end{document}
